\begin{theorem}[\textbf{Quantum Hamming Bound}]
    For a nondegenerate (pure) quantum code $[[n, k, d]]$ QEC code with alphabet size $q$ and $t=\lfloor\frac{d-1}{2}\rfloor$:
    \begin{equation}
        \sum_{j=0}^t\binom{n}{j} \left(q^2-1\right)^j\le \frac{q^n}{q^k}
    \end{equation}
    In the case of qubit, i.e., $q=2$, the bounds becomes:
    \begin{equation}
        \sum_{j=0}^t\binom{n}{j} 3^j\le 2^{n-k}
    \end{equation}
\end{theorem}

Unfortunately, the answer is no. 







After this step, within all the states that the Client receives, the Client can unambiguously determine $25\%$ of the states; even with the generalized POVM measurement, the Client can only unambiguously determine about $29\%$. For the states that the Client can unambiguously determine, he notes down the state; otherwise, he notes down randomly 0 or 1.

\begin{enumerate}
    \item The Client publicly announces several qubits that it can determine unambiguously.

    \item The Bank will check against the state he sends, and estimate the error during the channel. If the error is too high, they abort the protocol.

    \item The Bank will randomly
\end{enumerate}


    %Using Schmidt decomposition, any pure bi-partite quantum states $\ket{\psi}\in\mathcal{H}_A\otimes\mathcal{H}_B$ can be written as 
    %\begin{equation}
    %\ket{\psi}=\sum_{i=1}^{r}\lambda_i\ket{i_A}\otimes\ket{i_B}
    %\end{equation}
    %where $\ket{i_A}$ and $\ket{i_B}$ are local orthonormal basis, and non-negative Schmidt coefficients $\lambda_i$. Where the Schmidt rank $r\le \min(d_A, d_B)$ equals the number of nonzero Schmidt coefficients.


    Here, we would like to investigate the case that, suppose in Alice's lab, the measurement is asymmetric, which means, when Alice measures in $\{\ket{H}, \ket{V}\}$ basis, the error probability is $p_1$, which means, instead of measuring in $\{\ket{H}, \ket{V}\}$, due to the faulty apparatus, Alice is actually measuring in $\{\ket{H'}, \ket{V'}\}$, we assume that the error is small, so here we do not consider the phase information, where $\ket{H'}=\sqrt{1-p_1}\ket{H}+e^{i\theta}\sqrt{p_1}\ket{V}$ and $\ket{V'}=\sqrt{p_1}\ket{H}-\sqrt{1-p_1}\ket{V}$. Also when Alice measures in $\{\ket{D}, \ket{A}\}$ basis, the error probability is $p_2$, meaning that due to the faulty apparatus, Alice is actually measuring in $\{\ket{D'}, \ket{A'}\}$, where $\ket{D'}=\sqrt{1-p_2}\ket{D}+\sqrt{p_2}\ket{A}$ and $\ket{A'}=\sqrt{p_2}\ket{D}-\sqrt{1-p_2}\ket{A}$. that $p_1\ne p_2$. Without loss of generality, suppose $p_{thr}>p_1>p_2$, and $p_1-p_2=\delta_p$.

This is a realistic asymmetric error model, because the asymmetric error in different bases can be due to the quantum communication cable or the measurement device itself. Assume the overall error rate is $p_e=(p_1+p_2)/2$.

Bob might be able to cheat by shuffling the blocks differently. If Bob is not afraid to leak extra information to Alice, and Bob only wants to know which bit of information Alice wants, in each block, Bob can only put qubits prepared in $\{\ket{H}, \ket{V}\}$ basis or $\{\ket{D}, \ket{A}\}$ basis. When Alice chooses the measurement basis randomly, as we argued before, if Bob prepares the qubit in one basis, a unambiguous state discrimination is only possible when Alice measures in the complementary basis.

Thus, suppose a qubit is prepared in $\ket{H}$ and claim that the state is either $\ket{H}$ or $\ket{A}$, and Alice measures it in the $\{\ket{D'}, \ket{A'}\}$, where the probability for a USD is 
\begin{equation}
\begin{split}
    |\bra{H}\ket{D'}|^2&=\lvert\sqrt{1-p_2}\bra{H}\ket{D}+e^{i\theta}\sqrt{p_2}\bra{H}\ket{A}\rvert^2\\
    &=p \sin ^2(\theta)+\frac{1}{2} \left(\sqrt{p} \cos (\theta)+\sqrt{1-p}\right)^2\\
\end{split}
\end{equation}


The first observation we have is in fact that, if Alice wants to execute this kind of cheating, then she must do measurements as normal, and only during the channel error estimation stage, she lies about the measurement result she gets. Therefore, a way to prevent this is to catch Alice lying. This can be done with extra resources, utilizing the CHSH inequality.

We would like to quickly recall that in the CHSH game, the 2 parties have a shared Bell state $\ket{\Phi^+}=\frac{1}{\sqrt{2}}(\ket{00}+\ket{11})$, and one party randomly chooses a basis from $A_0=\sigma_X, A_1=\sigma_Z$ to measure. While the other party chooses randomly from $B_0=\frac{\sigma_X+\sigma_Z}{\sqrt{2}}, B_1=\frac{\sigma_X-\sigma_Z}{\sqrt{2}}$.

The eigenstates of the measurement operators are listed below, where $c=\cos\left(\frac{\pi}{8}\right)$ and $s = \sin\left(\frac{\pi}{8}\right)$
    \begin{equation}
        \ket{A_0^+}=\ket{+}=\frac{\ket{0}+\ket{1}}{\sqrt{2}}, \ket{A_0^-}=\ket{-}=\frac{\ket{0}-\ket{1}}{\sqrt{2}}
    \end{equation}
    \begin{equation}
        \ket{A_1^+}=\ket{0}, \ket{A_1^-}=\ket{1}
    \end{equation}

    \begin{equation}
    \begin{split}
        \ket{B_0^+}&=c\ket{0}+s\ket{1}=c\ket{+}+s\ket{-},\\
        \ket{B_0^-}&=s\ket{0}-c\ket{1} =s\ket{+}-c\ket{-}
    \end{split}
    \end{equation}
    
    \begin{equation}
    \begin{split}
        \ket{B_1^+}&=c\ket{0}-s\ket{1}=s\ket{+}+c\ket{-},\\
        \ket{B_1^-}&=s\ket{0}+c\ket{1} = c\ket{+}-s\ket{-}
    \end{split}
    \end{equation}

We modify the first few steps of the protocol in Section~\ref{sec:the_protocol} as follows: 
\begin{enumerate}
    \item When Bob sends the states to Alice, he will also send some states prepared in $\{\ket{B_0^+}, \ket{B_0^-}\}$ or $\{\ket{B_1^+}, \ket{B_1^-}\}$. For easy reference, we called the states prepared in $\ket{0}, \ket{1}, \ket{+}, \ket{-}$ as the ``normal'' states, and $\ket{B_0^+}, \ket{B_0^-},\ket{B_1^+}, \ket{B_1^-}$ as the ``twisted'' states.
    
    \item \textbf{Channel error estimation stage}: They first choose several normal states to publicly announce the measurement basis and result, and get an estimation of the channel error rate. But remember that, Alice are intended to lie about several measurement results to exaggerate a higher channel error rate.

    \item \textbf{Catch Alice lying stage}: Then, Bob will randomly select several extra states, containing both normal and twisted states, and ask Alice to publicly announce the measurement basis and results, while he is not announcing. 

    \item Since Alice do not know which states are normal and which states are twisted, the best she can do is to use the same cheating strategy as before, randomly faking some measurement results.

    \item When Bob gets the Alice's announcement, he first check for all normal states he selected, that the channel error rate appears to be agreeing with the value they get from the channel error estimation stage.

    \item Then, Bob checks for the twisted states he selected. We will demonstrate how Bob can catch Alice cheating.
\end{enumerate}

So, if Bob sends the state $\ket{B_0^+}=c\ket{0}+s\ket{1}=c\ket{+}+s\ket{-}$, no matter what basis did Alice measures, she should get the $+1$ eigenvalue with probability $p=\cos^2\left(\frac{\pi}{8}\right)\approx 0.85$. 



Firstly, we define the Hamming ball as follows,

\begin{definition}
    For integers $q, n, \ell$, the Hamming ball of radius $\ell$ in $\{0, 1, \cdots, q-1\}$ is denoted by $\mathrm{Vol}_q(n, \ell)$, and it is given by:
    \begin{equation}
        \mathrm{Vol}_q(n, \ell) = \sum_{j=0}^{\ell} \binom{n}{j} (q - 1)^j.
    \end{equation}
\end{definition}